\documentclass{article}
\usepackage[utf8]{inputenc}
\usepackage[T1]{fontenc}
\usepackage[a4paper]{geometry}		% Reduire les marges
\geometry{hmargin=3cm,vmargin=3.5cm}
\usepackage[english]{babel} 
\usepackage{amsfonts}
\usepackage{amsthm}
\usepackage{amsmath}
\usepackage{amssymb}
\usepackage{hyperref}
\usepackage{array}
\usepackage[svgnames]{xcolor}
\usepackage{xpatch}
\usepackage{tikz}
\usepackage{mathrsfs}
%\usepackage{graphicx}
%\usepackage{caption}
\usepackage[svgnames]{xcolor}
\usepackage{eurosym}
\usetikzlibrary{positioning,calc}
\renewcommand{\thesection}{\Roman{section}}
\xpatchcmd{\proof}{\itshape}{\normalfont\proofnameformat}{}{}
\newcommand{\proofnameformat}{\bfseries}
\theoremstyle{definition}
\newtheorem{defi}{Definition}[section]
\newtheorem{theo}{Theorem}[section]
\newtheorem{prop}[theo]{Proposition}
\newtheorem{lem}[theo]{Lemma}
\newtheorem{coro}[theo]{Corollary}
\newtheorem{exe}{Example}[section]
\newtheorem{rem}{Remark}[section]
\newtheorem{nota}{Notation}[section]
\renewcommand\qedsymbol{$\blacksquare$}
\numberwithin{equation}{section}

\DeclareMathOperator{\R}{\mathbb{R}}
\DeclareMathOperator{\Z}{\mathbb{Z}}
\DeclareMathOperator{\Q}{\mathbb{Q}}
\DeclareMathOperator{\N}{\mathbb{N}}
\DeclareMathOperator{\E}{\mathbb{E}}
\DeclareMathOperator{\C}{\mathcal{C}}

\title{Exo 10, corrigé}       
\author{Guillaume Woessner}
\date{}                       % sans date : celle du jour de compilation

%\pagestyle{headings}          % Pour mettre des entetes avec les titres
                              % des sections en haut de page
\sloppy                       % Ne pas faire deborder les lignes dans la marge

\begin{document}

\maketitle                    % Faire un titre utilisant les donnees
                              % passees : \title, \author et \date

\begin{abstract}
\begin{center}
Le mouvement brownien (MB).
\end{center}
\end{abstract}

%\tableofcontents              % Table des matieres


Dans tout ce document, l'espace de probabilité filtré considéré est $(\Omega,\mathcal{F},(\mathcal{F}_n)_{t\in[0,T]},\mathbb{P})$, et on supposera que $(B_t)_t$ est un mouvement brownien adapté à la filtration.


\paragraph{Exercice 1}
Dans le cours on a construit le mouvement brownien sur $[0,1]$. Dans cet exercice on va l'étendre à $[0,\infty[$.\\
Soient, pour $i=1,2,\dots$, des MB indépendants $(B_t^i)_{t\in[0,1]}$. On définit alors le processus $(B_t)_{t\geq 0}$ par 
$$ B_t := \sum_{i=1}^{\lfloor t \rfloor} B^i_1 + B^{\lfloor t \rfloor + 1}_{t-\lfloor t \rfloor}. $$
Montrez que le processus $(B_t)_{t\geq 0}$ est bien un mouvement brownien sur $[0,\infty[$.

\begin{proof}
On va vérifier que les 4 propriétés du MB sont bien vérifiées.
\begin{itemize}
    \item[(1)] On a bien évidemment $B_0=0$, car $\sum_{i=1}^0=0$ par définition.
    \item[(2)] Remarquons d'abord que, pour $s<t$, 
    \begin{align}
    \label{rouge}
    B_t-B_s = \sum_{i=\lfloor s \rfloor + 1}^{\lfloor t \rfloor} B^i_1 + B_{t-\lfloor t \rfloor}^{\lfloor t \rfloor + 1} - B_{s-\lfloor s \rfloor}^{\lfloor s \rfloor + 1}
    \end{align}
    Ensuite, pour $t_1 \leq t_2 \leq t_3$, il s'agit de montrer que $(B_{t_3}-B_{t_2}) \amalg (B_{t_2}-B_{t_1})$. On peut procéder par disjonction de 4 cas, selon que $\lfloor t_1\rfloor < \lfloor t_2\rfloor < \lfloor t_3\rfloor$, ou que l'une ou les deux inégalités soient des égalités. On présente ici le premier cas.\\
    Dans ce cas on a immédiatement que $\sum_{i=\lfloor t_1 \rfloor+1}^{\lfloor t_2\rfloor} B_1^i$ et $-B^{\lfloor t_1\rfloor+1}_{t_1-\lfloor t_1\rfloor}$ sont indépendants de $B_{t_3}-B_{t_2}$.\\
    Le terme restant $B_{t_2-\lfloor t_2\rfloor}^{\lfloor t_2\rfloor+1}$ est lui indépendant de $B_1^{\lfloor t_2\rfloor+1}-B_{t_2-\lfloor t_2\rfloor}^{\lfloor t_2\rfloor +1}$, ainsi que de tous les autres termes apparaissant dans $B_{t_3}-B_{t_2}$.\\
    On a donc bien l'indépendance recherchée dans le premier cas, et je vous laisse faire les autres vous-même............
    \item[(3)] On veut montrer que pour $s\leq t$, on a bien $B_t-B_s \sim \mathcal{N}(0,t-s)$. Or en réécrivant \eqref{rouge} on obtient, si $\lfloor s \rfloor < \lfloor t \rfloor$ (l'autre cas est à nouveau laissé à l'initiative du lecteur),
    $$ B_t-B_s = (B^{\lfloor s \rfloor +1}_1-B^{\lfloor s \rfloor +1}_{s-\lfloor s \rfloor}) + \sum_{i=\lfloor s \rfloor + 2}^{\lfloor t \rfloor} B^i_1 + B_{t-\lfloor t \rfloor}^{\lfloor t \rfloor + 1}, $$
    ce qui correspond à une somme de gaussiennes indépendantes
    $$ \mathcal{N}(0,1-(s-\lfloor s \rfloor)) + \sum_{i=\lfloor s \rfloor + 2}^{\lfloor t \rfloor} \mathcal{N}(0,1) + \mathcal{N}(0,t-\lfloor t \rfloor) ~~\sim~~ \mathcal{N}(0,t-s). $$
    \item[(4)] Chaque MB $B^i$ est continu sur un ensemble presque sûr $\Omega_i\subset \Omega$, ainsi $B$ est bien continu sur $\cap_i \Omega_i$, qui est lui aussi presque sûr.
\end{itemize}
\end{proof}



\paragraph{Exercice 2}
Soit $X:=(X_1,\dots,X_n)$ un vecteur gaussien. Montrez qu'il est caractérisé par son vecteur d'espérance $m_i:=\E[X_i]$ et sa matrice de covariance $Cov(i,j):=\E[X_iX_j]$.

\begin{proof}
On calcule la fonction caractéristique du vecteur $X$, donnée par
$$\varphi_X(u):=\E\left[e^{i\langle u~|~X\rangle}\right] = \exp\left( i\E[\langle u~|~X\rangle] - \dfrac{1}{2}Var(\langle u~|~X\rangle) \right).$$
Or, $\E[\langle u~|~X\rangle] = \langle u~|~\E[X]\rangle = \langle u~|~m\rangle$, et 
\begin{align*}
Var(\langle u~|~X\rangle) &= \E\left[ (\langle u~|~X\rangle)^2 \right]-\E[(\langle u~|~m\rangle)^2] \\
&= \sum_{ij}u_iu_j\E[X_iX_j]=u^t.Cov.u.
\end{align*}
Ainsi la fonction caractéristique, et donc la loi, du vecteur gaussien $X$ est bien caractérisée par espérance et matrice de covariance.
\end{proof}

\paragraph{Exercice 3}
Trouver un processus $(W_t)_{t\in[0,1]}$, qui a les mêmes lois marginales que $(B_t)_{t\in[0,1]}$ mais qui n'est pas continu.

\begin{proof}
On peut prendre $W_t := B_t + \delta_Y(t)$ avec $Y\sim\mathcal{U}[0,1]$ indépendante, comme dans la Feuille 9 Exercice 4.
\end{proof}

\paragraph{Exercice 4}
Calculer $\E[B_5~|~B_1]$ et $\E[B_5^2~|~B_1]$.

\begin{proof}
On a $$\E[B_5~|~B_1]=\E[B_5-B_1~|~B_1]+B_1=B_1,$$
et
$$\E[B_5^2~|~B_1]=\E[(B_5-B_1)^2+2B_5B_1-B_1^2~|~B_1]=\E[(B_5-B_1)^2]+2B_1^2-B_1^2=4+B_1^2.$$
\end{proof}


\paragraph{Exercice 5}
Fixons $t>0$ et posons, pour $n\in\N$ et $0\leq i < 2^n$ des entiers, $t_i:=t\frac{i}{2^n}$.\\
Montrez que
$$\sum_{i=0}^{2^n-1} (B_{t_{i+1}}-B_{t_i})^2 ~\mapsto t,$$
dans $L^2(\mathbb{P})$ et \textit{presque sûrement}.

\begin{proof}
Commençons par la preuve de la convergence $L^2$.\\
On note, pour tout $n$ et $i$, $Y^n_i:=(B_{t_{i+1}}-B_{t_i})^2-\frac{t}{2^n}$. A $n$ fixé les $(Y_i^n)_i$ sont indépendantes et centrées, et vérifient $Var(Y_i^n)=\E[(Y_i^n)^2]=^\star2(\frac{t}{2^n})^2$. Ainsi on calcule,
\begin{align*}
\E\left[\left( \sum_{i=0}^{2^n-1} (B_{t_{i+1}}-B_{t_i})^2-t \right)^2\right] &= \E\left[ \left( \sum_{i=0}^{2^n-1}Y_i^n \right)^2 \right] = \sum_{i=0}^{2^n-1}Var(Y^n_i)\\
&=2\sum_{i=0}^{2^n-1} \left(\dfrac{t}{2^n}\right)^2 = \dfrac{2}{2^n}t^2 ~~ \longrightarrow_n~~ 0.
\end{align*}
Passons maintenant à la convergence presque sûre.\\
On a vu que 
$$ \E\left[ \left( \sum_{i=0}^{2^n-1}Y_i^n \right)^2 \right] = \dfrac{t^2}{2^{n-1}}. $$
Ainsi, par l'inégalité de Chebychev
$$ \mathbb{P}\left( \left| \sum_{i=0}^{2^n-1} Y_i^n \right| > \dfrac{1}{n} \right) \leq \dfrac{t^2n^2}{2^{n-1}},$$
qui est le terme générale d'une série sommable en $n\geq 1$. Ainsi d'après le lemme de Borel-Cantelli, il existe presque sûrement un rang $N$ à partir duquel pour $n\geq N$ on ait
$$  \left| \sum_{i=0}^{2^n-1} Y_i^n \right| \leq \dfrac{1}{n},$$
ce qui implique bien la convergence presque sûre.
\end{proof}


\paragraph{Exercice 8}
Calculer $\mathbb{P}(B_1>0 \cap B_2>0)$ et $\mathbb{P}(B_1>0 \cap B_2<0)$.

\begin{proof}
On calcule
\begin{align*}
    \mathbb{P}(B_1>0 \cap B_2>0) &= \int_{\R} \mathbb{P}(B_2>0\cap B_1>0~|~B_1=
    y)d\mathbb{P}(B_1=y)\\
    &=\int_{\R^+}\dfrac{1}{\sqrt{2\pi}}e^{-\frac{y^2}{2}}\mathbb{P}(B_2>0~|~B_1=y)dy\\
    &=\int_{\R^+}\dfrac{1}{\sqrt{2\pi}}e^{-\frac{y^2}{2}}\mathbb{P}(B_2-B_1>-y~|~B_1=y)dy\\
    &=\int_{\R^+}\dfrac{1}{\sqrt{2\pi}}e^{-\frac{y^2}{2}}\mathbb{P}(B_2-B_1>-y)dy\\
    &=\int_{\R^+}\dfrac{1}{\sqrt{2\pi}}e^{-\frac{y^2}{2}}\int_{-y}^\infty \dfrac{1}{\sqrt{2\pi}}e^{-\frac{z^2}{2}}dzdy\\
    &=\int_{\R^+}\dfrac{1}{\sqrt{2\pi}}e^{-\frac{y^2}{2}}\int_0^\infty \dfrac{1}{\sqrt{2\pi}}e^{-\frac{z^2}{2}}dzdy + \int_{\R^+}\dfrac{1}{\sqrt{2\pi}}e^{-\frac{y^2}{2}}\int_{-y}^0 \dfrac{1}{\sqrt{2\pi}}e^{-\frac{z^2}{2}}dzdy\\
    &=\int_{\R^+}\dfrac{1}{\sqrt{2\pi}}e^{-\frac{y^2}{2}}dy\int_0^\infty \dfrac{1}{\sqrt{2\pi}}e^{-\frac{z^2}{2}}dz + \int_{\R^+}\dfrac{1}{\sqrt{2\pi}}e^{-\frac{y^2}{2}}\int_0^y \dfrac{1}{\sqrt{2\pi}}e^{-\frac{z^2}{2}}dzdy\\
    &=\dfrac{1}{2}.\dfrac{1}{2}+\int_{\R^+}\dfrac{1}{\sqrt{2\pi}}e^{-\frac{y^2}{2}}.\left(\dfrac{1}{2}\text{erf}(\sqrt{2}y)\right)dy\\
    &=\dfrac{1}{4}+\int_{\R^+}\dfrac{1}{2\sqrt{\pi}}e^{-\frac{x^2}{4}}\text{erf}(x)dx,
\end{align*}
où erf est la fonction d'erreur de Gauss.\\
Pour le deuxième calcul, il faut utiliser que 
$$ \mathbb{P}(B_1>0\cap B_2>0)+\mathbb{P}(B_1>0\cap B_2<0)= \mathbb{P}(B_1>0)=\dfrac{1}{2}. $$
\end{proof}

~\\
$\star$ il faut utiliser que si $B\sim \mathcal{N}(0,\sigma^2)$, alors $\E[B^{2k}]=\frac{(2k)!}{2^kk!}\sigma^{2k}$.

\end{document}


\paragraph{Exercice 6}
Exercice pour montrer les propriétés de symétrie, d'échelle, d'inversion du temps et de retournement du temps ?


\paragraph{Exercice 8*}
Soit $(\varphi_n)_n$ une base hilbertienne de $H:=L^2([0,T],\mathcal{B},dx)$, et $(\epsilon_n)_n$ une famille \textit{i.i.d.} de variables $\mathcal{N}(0,1)$.\\
Montrez que la famille de processus $(B_t^n)_n$ définie par 
$$ B_t^n := \sum_{k=0}^n \epsilon_k \langle \varphi_k~|~\mathbf{1}_{[0,t]}\rangle,$$
où $\langle\cdot~|~\cdot\rangle$ est le produit scalaire de $H$, converge dans $L^2(\mathbb{P})$ vers un mouvement brownien.


\paragraph{Exercice 6}
Soit $\lambda\geq 0$.\\
Montrez que les processus $(M_t)_{t\in[0,1]}$ défini par $M_t := e^{\lambda B_t - \frac{\lambda^2}{2}t}$ est une martingale.